\section{evnx cloud}

\begin{frame}{\cmd{evnx cloud} — encrypted sync}
  \begin{block}{Synopsis}\texttt{evnx cloud <SUBCOMMAND>}\end{block}

  \begin{tabular}{@{}ll@{}}
    \toprule
    \texttt{cloud push}   & Encrypt locally, upload ciphertext \\
    \texttt{cloud pull}   & Download and decrypt \\
    \texttt{cloud run -- <CMD>} & Inject into a subprocess; nothing on disk \\
    \texttt{cloud status} & Server, account, token state \\
    \texttt{cloud history}& Version list \\
    \texttt{cloud link <VAULT>} & Bind this directory to a vault \\
    \texttt{cloud unlink} & Remove the binding \\
    \bottomrule
  \end{tabular}
\end{frame}

\begin{frame}[fragile]{The zero-knowledge chain}
\begin{semiverbatim}\scriptsize
Password
  -> Argon2id(password, salt, m=64MB, t=3, p=4)
       -> MasterKey            (never leaves your device)
           -> XChaCha20(MK) wraps VaultKey   (random 256-bit, one per vault)
               -> AES-256-GCM(VK, nonce) encrypts .env
                   -> ciphertext -> object storage
\end{semiverbatim}

  \vspace{0.5em}
  \begin{block}{Associated data stops a rollback}
    Blobs are encrypted with \texttt{vault\_aad(vault\_id, version)} bound in, so
    a malicious server cannot replay an old version as current.

    \vspace{0.3em}
    {\footnotesize Consequence: a \texttt{409} conflict requires
    \textbf{re-encryption}, not a bare retry — the version number is
    authenticated into the ciphertext.}
  \end{block}
\end{frame}

\begin{frame}[fragile]{\cmd{cloud run} — secrets without a file}
  \capturepart{61-cloud-run-help.txt}{1}{12}
\end{frame}

\begin{frame}{Why \cmd{cloud run} is different}
  \begin{itemize}
    \item Values travel in the \textbf{environment block}, never the argument
          vector — so nothing in \texttt{ps}, nothing in shell history,
          nothing on disk.
    \item \flag{--include} / \flag{--exclude} narrow what the child sees.
          A build step needing only \texttt{NEXT\_PUBLIC\_*} has no reason to
          hold your database password.
    \item A filter matching nothing is an \textbf{error}, not a silent run with
          zero secrets.
    \item The child's exit code comes back unchanged.
    \item It runs the command directly, not through a shell. For a pipeline,
          ask for one: \texttt{-- sh -c 'a | b'}.
  \end{itemize}

  \vspace{0.6em}
  \begin{alertblock}{Be precise about what it removes}
    It removes the plaintext \emph{file}, not the master password. The vault key
    is wrapped under your master key, so decryption needs it wherever this runs.
    In CI, feed it with \flag{--password-stdin}.
  \end{alertblock}
\end{frame}

\begin{frame}[fragile]{\cmd{cloud status}}
  \capture{65-cloud-status.txt}

  {\footnotesize A directory bound with \cmd{cloud link} needs no
  \flag{--vault} on any subsequent command.}
\end{frame}
